
dispose de l'anneau de Chow @@M0@@ des classes de cycles à équivalence rationnelle près [4]. Comme variantes de cet anneau, défini d'ailleurs sous des considérations sensiblement plus générales, on a les anneaux @@M1@@ (gradué associé à la @@M2@@-filtration de @@M3@@) et @@M4@@ (gradué associé à la filtration de @@M5@@ par la codimension des supports). On a des théories des classes de Chern à valeurs dans chacun de ces trois anneaux (cf. [9] pour le cas de @@M6@@, et ce séminaire pour @@M7@@, qui s'envoie dans @@M8@@). On a également une théorie cohomologique (@@M9@@-adique, disons, où @@M10@@ est premier à la caractéristique de @@M11@@) des classes de Chern (SGA 5 VII)~:
@@M12@@
(où @@M13@@ est un fibré vectoriel sur @@M14@@, ou encore un élément quelconque de @@M15@@). Ces théories cohomologiques sont reliées à l'aide des homomorphismes canoniques (compatibles avec les structures multiplicatives)
@@M16@@
où on a posé
@@M17@@
L'homomorphisme @@M18@@ provient du fait que la @@M19@@-filtration de @@M20@@ est plus fine que la filtration topologique par la codimension du support (Exp. X), et c'est un isomorphisme modulo torsion (Exp. VII 4.11). L'homomorphisme @@M21@@ est défini par exemple dans [RRR], et associe à la classe du cycle de codimension @@M22@@ associé à un sous-schéma fermé @@M23@@ de @@M24@@ de codimension @@M25@@, l'élément @@M26@@ de @@M27@@ mod @@M28@@. C'est un homomorphisme surjectif, dont nous allons voir encore (4.2) que c'est un isomorphisme modulo torsion. Ainsi, \uline{modulo torsion les trois anneaux} @@M29@@, @@M30@@ \uline{et} @@M31@@ \uline{sont canoniquement isomorphes}, donc \uline{à posteriori} les formules de Riemann--Roch exprimées à l'aide de l'un ou l'autre de ces anneaux (tensorisés par @@M32@@) sont équivalentes. L'anneau gradué @@M33@@, surtout utilisé dans ce séminaire comme substitut à l'anneau de Chow, a l'avantage appréciable d'être défini pour tout topos localement annelé, et en particulier pour tout schéma 
